No meio do software livre, os textos (códigos-fonte) que compõem os diferentes softwares são abertamente publicados. Muitas vezes, passam a sofrer intervenções pelos usuários, que republicam o código modificado num processo contínuo e coletivo de transformação do software. Proponho a carta de intenções como uma publicação científica cujas notas de rodapé sejam passíveis de reescritas, localizando o texto a cada ocasião que assim se queira. Cada reescrita configura uma iteração do texto, sendo a avaliação da carta realizável como uma iteração sua.